\section{Proof of the general formula and attainment}\label{sec:main-proof}

We retain the notation of Theorem~\ref{thm:main} and the block covariance in \eqref{eq:full-covariance}.
All traces in this section are unnormalized matrix traces.
Adjoints of maps are taken for the Hilbert--Schmidt inner product unless stated otherwise.

\subsection{A coercive regularization}
For $\varepsilon>0$ and $c>0$, put
\[
 C:=\varepsilon^21_m+\eta(c),\qquad
 \mathcal L_{\varepsilon,c}(Y):=CYc^{-1}+\rho(Y^*).
\]
On the real vector space $M_m(\C)$ use the inner product $\langle X,Y\rangle_{\R}:=\operatorname{Re}\Tr(X^*Y)$.
The map $\mathcal L_{\varepsilon,c}$ is selfadjoint for this inner product, and its quadratic form is
\begin{equation}\label{eq:quadratic-form}
 B_{\varepsilon,c}(Y):=\Tr(Y^*CYc^{-1})
             +\operatorname{Re}\Tr(Y^*\rho(Y^*)).
\end{equation}
For the second summand, symmetry follows from cyclicity of the trace:
$\operatorname{Re}\Tr(X^*a_jY^*a_j) =\operatorname{Re}\Tr(Y^*a_jX^*a_j)$.
To check positivity, set $D_j:=c^{-1/2}Y^*a_jc^{1/2}$.
Then
\[
 |\Tr(D_j^2)|\le\Tr(D_j^*D_j)
      =\Tr(Y^*a_jca_j^*Yc^{-1}).
\]
The inequality is the Hilbert--Schmidt Cauchy--Schwarz inequality applied to $D_j^*$ and $D_j$.
Consequently,
\begin{equation}\label{eq:quadratic-lower-bound}
 B_{\varepsilon,c}(Y)\ge
           \varepsilon^2\Tr(Y^*Yc^{-1})>0\qquad(Y\ne0).
\end{equation}
Thus $\mathcal L_{\varepsilon,c}$ is invertible as a real-linear map on $M_m(\C)$.
Its unique solution $Y=Y_{\varepsilon,c}$ at right-hand side $a_0$ depends smoothly on $c$.
In particular, $\varepsilon=1$, $c=b$ proves the asserted invertibility of $\mathcal T_b$.

Define
\begin{equation}\label{eq:regularized-objective}
 \mathcal F_\varepsilon(c):=\log\det C-\log\det c
       +\Tr c-m+\operatorname{Re}\Tr(a_0^*Y_{\varepsilon,c}).
\end{equation}
Completing the square in the real inner product gives
\begin{equation}\label{eq:inner-maximum}
 \operatorname{Re}\Tr(a_0^*Y_{\varepsilon,c})
 =\max_Y\{2\operatorname{Re}\Tr(a_0^*Y)-B_{\varepsilon,c}(Y)\}
 \ge0.
\end{equation}
Hence
\[
 \mathcal F_\varepsilon(c)\ge
 m\log\varepsilon^2+\Tr c-\log\det c-m.
\]
Since $t-\log t\to+\infty$ as $t\downarrow0$ or $t\to\infty$, every sublevel set is contained in $\{\alpha 1_m\le c\le\beta 1_m\}$ for some $0<\alpha\le\beta<\infty$.
Continuity then shows that $\mathcal F_\varepsilon$ attains a global minimum in the positive cone.

\subsection{Critical points and the resolvent solution}
We now identify every critical point of $\mathcal F_\varepsilon$.
In this subsection $Y=Y_{\varepsilon,c}$.
Differentiating the maximum in \eqref{eq:inner-maximum} requires only the explicit $c$-derivative: the derivative with respect to $Y$ vanishes at its maximizer.
Put
\begin{equation}\label{eq:auxiliary-P}
 P:=C^{-1}-Yc^{-1}Y^*.
\end{equation}
This matrix is initially only known to be selfadjoint.
Differentiating \eqref{eq:regularized-objective} in every selfadjoint direction gives
\begin{equation}\label{eq:regularized-critical}
 c^{-1}=1_m+\eta^*(P)+c^{-1}Y^*CYc^{-1}.
\end{equation}
Indeed, the terms involving the derivative of $C$ combine to $\Tr(\eta^*(P)\,dc)$; the derivative of $c^{-1}$ in the quadratic form contributes $\Tr(c^{-1}Y^*CYc^{-1}\,dc)$.

We next prove that $P$ is positive definite.
Set
\[
 U:=C^{1/2}Yc^{-1/2},\qquad K:=1_m-U^*U,\qquad W:=1_m-UU^*,
\]
and define the completely positive map
\[
 \Phi(x):=C^{-1/2}\eta(c^{1/2}xc^{1/2})C^{-1/2}.
\]
Then $P=C^{-1/2}WC^{-1/2}$, and \eqref{eq:regularized-critical} is equivalent to
\begin{equation}\label{eq:sign-fixed-point}
 K=c+\Phi^*(W),\qquad
 \Phi(1_m)=1_m-\varepsilon^2C^{-1}<1_m.
\end{equation}
Choose a unitary extension $V$ of the polar factor of the square matrix $U$, so that $U=V(U^*U)^{1/2}$ and $W=VKV^*$.
For $\Lambda(x):=\Phi^*(VxV^*)$ we have $\Lambda^*(1_m)=V^*\Phi(1_m)V<1_m$.
The positive-map estimate in Section~\ref{subsec:gap-tools} implies that the spectral radius of $\Lambda$, equal to that of its adjoint, is less than one.
Thus
\[
 K=(\id-\Lambda)^{-1}(c)=\sum_{k\ge0}\Lambda^k(c)>0.
\]
It follows that $W>0$ and $P>0$ as required.

Consider the invertible selfadjoint block matrix
\[
 \mathfrak q:=\begin{pmatrix}P&Y\\Y^*&-c\end{pmatrix}.
\]
Its Schur complement is $P+Yc^{-1}Y^*=C^{-1}$, so
\begin{equation}\label{eq:signed-inverse}
 \mathfrak q^{-1}=
 \begin{pmatrix}
 C&CYc^{-1}\\
 c^{-1}Y^*C&-c^{-1}+c^{-1}Y^*CYc^{-1}
 \end{pmatrix}
 =\begin{pmatrix}\varepsilon^21_m&a_0\\a_0^*&-1_m\end{pmatrix}
       -\Theta(\mathfrak q).
\end{equation}
The upper right equality uses $CYc^{-1}+\rho(Y^*)=a_0$, and the lower diagonal equality is \eqref{eq:regularized-critical}.
Now set
\begin{equation}\label{eq:physical-block}
 \mathscr R:=\begin{pmatrix}i\varepsilon P&Y\\Y^*&ic/\varepsilon\end{pmatrix}.
\end{equation}
Multiplying \eqref{eq:signed-inverse} by the corresponding diagonal scalars, or directly inverting \eqref{eq:physical-block}, gives
\[
 \mathscr R^{-1}=A-\Theta(\mathscr R)-i\varepsilon 1_{2m}.
\]
Since $\operatorname{Im}\mathscr R>0$, the matrix $-\mathscr R$ is the unique solution of \eqref{eq:dyson-general} at $i\varepsilon$.
Thus $\mathscr R=E_{2m}((\mathscr H-i\varepsilon 1_{2m})^{-1})$.
The operator block inverse yields
\begin{equation}\label{eq:regularized-optimizer}
 \begin{split}
 P_\varepsilon&=E_m((SS^*+\varepsilon^21_m)^{-1}),\\
 Y_\varepsilon&=E_m((SS^*+\varepsilon^21_m)^{-1}S),\\
 c_\varepsilon&=\varepsilon^2E_m((S^*S+\varepsilon^21_m)^{-1}).
 \end{split}
\end{equation}
Every critical point therefore has the same $c=c_\varepsilon$.
Since a minimum exists, $c_\varepsilon$ is the unique minimizer and the unique critical point of $\mathcal F_\varepsilon$.

\subsection{The minimum value and removal of regularization}
All quantities in the next calculation are evaluated at \eqref{eq:regularized-optimizer}.
Write $\gamma_\rho:=\operatorname{Re}\Tr(Y^*\rho(Y^*))$.
The Schur complement gives $|\det\mathscr R|=\det c/\det C$, while block multiplication gives
\[
 \tfrac12\operatorname{Re}\Tr(\Theta(\mathscr R)\mathscr R)
       =-\Tr(P\eta(c))+\gamma_\rho.
\]
Applying \eqref{eq:regularized-primitive} to $G_\varepsilon=-\mathscr R$ therefore gives
\begin{equation}\label{eq:minimum-primitive}
 2m\log\Delta_\varepsilon(S)
 =\log\det C-\log\det c-\Tr(P\eta(c))+\gamma_\rho.
\end{equation}
On the other hand, $\mathcal L_{\varepsilon,c}(Y)=a_0$ and the trace of \eqref{eq:regularized-critical} after multiplication by $c$ give
\[
 \operatorname{Re}\Tr(a_0^*Y)=\Tr(Y^*CYc^{-1})+\gamma_\rho,
 \qquad
 m=\Tr c+\Tr(P\eta(c))+\Tr(Y^*CYc^{-1}).
\]
Substitution in \eqref{eq:minimum-primitive} proves
\begin{equation}\label{eq:regularized-minimum}
 2m\log\Delta_\varepsilon(S)=\min_{c>0}\mathcal F_\varepsilon(c).
\end{equation}

The substitution $c=\varepsilon^2b$ is a bijection of the positive cone.
It gives $C=\varepsilon^2H_b$ and $\mathcal L_{\varepsilon,c}=\mathcal T_b$, hence
\begin{equation}\label{eq:scaled-minimum}
 2m\log\Delta_\varepsilon(S)
      =\min_{b>0}\{\mathcal V(b)+\varepsilon^2\Tr b\}.
\end{equation}
Let $\alpha:=\inf_{b>0}\mathcal V(b)\in[-\infty,\infty)$.
The right-hand side is at least $\alpha$.
For each fixed $b>0$ it is at most $\mathcal V(b)+\varepsilon^2\Tr b$, so its upper limit as $\varepsilon\downarrow0$ is at most $\mathcal V(b)$.
Taking the infimum over fixed trials proves convergence to $\alpha$, also when $\alpha=-\infty$.
Since the left-hand side decreases to $2m\log\Delta(S)$, we obtain $2m\log\Delta(S)=\inf_{b>0}\mathcal V(b)$.
Exponentiating proves \eqref{eq:main-determinant}.

\subsection{A critical point forces a spectral gap}
Suppose now that $b>0$ is a critical point of $\mathcal V$.
Write $H:=1_m+\eta(b)$, $Y:=Z_b$, and $P:=H^{-1}-Yb^{-1}Y^*$.
The same differentiation as above, with the term $\Tr c$ absent, gives
\begin{equation}\label{eq:unregularized-critical}
 b^{-1}=\eta^*(P)+b^{-1}Y^*HYb^{-1}.
\end{equation}
Set $U:=H^{1/2}Yb^{-1/2}$ and use the same $K,W,V,\Phi$ as in \eqref{eq:sign-fixed-point}, with $C=H$ and $c=b$.
This time $K=\Phi^*(W)=\Lambda(K)$ and $\Phi(1_m)=1_m-H^{-1}<1_m$.
Thus $K=0$.
Consequently $U$ is unitary, $P=0$, and $Y$ is invertible.
We have proved
\begin{equation}\label{eq:boundary-identities}
 b=Y^*(1_m+\eta(b))Y,\qquad
 a_0=Y^{-*}+\rho(Y^*),
\end{equation}
where $Y^{-*}:=(Y^*)^{-1}$.

These identities also give a spectral-radius bound on the full block space.
Define the completely positive maps
\[
 \mathcal A(x):=Y^*\eta(x)Y,\qquad
 \mathcal B(x):=Y\eta^*(x)Y^*.
\]
The first identity in \eqref{eq:boundary-identities} says $b-\mathcal A(b)=Y^*Y>0$, so the spectral radius of $\mathcal A$ is less than one.
The map $\mathcal B$ is similar to $\mathcal A^*:x\mapsto\eta^*(YxY^*)$, by conjugation with the invertible map $x\mapsto YxY^*$.
It too has spectral radius less than one.
Hence
\[
 P_0:=(\id-\mathcal B)^{-1}(YY^*)>0,
 \qquad P_0-\mathcal B(P_0)=YY^*.
\]
Put
\begin{equation}\label{eq:gap-basepoint}
 G_0:=-\begin{pmatrix}0&Y\\Y^*&0\end{pmatrix},\qquad
 D:=\begin{pmatrix}P_0&0\\0&b\end{pmatrix}>0,
 \qquad \mathcal S(x):=G_0\Theta(x)G_0.
\end{equation}
The second identity in \eqref{eq:boundary-identities} and its adjoint give $G_0^{-1}=-A-\Theta(G_0)$.
Moreover,
\begin{equation}\label{eq:strict-full-stability}
 D-\mathcal S(D)=G_0^2
       =\operatorname{diag}(YY^*,Y^*Y)>0.
\end{equation}
The positive-map estimate therefore gives $r(\mathcal S)<1$ on all of $M_{2m}(\C)$.

We explain why the selfadjoint solution $G_0$ is the value at zero of the analytically continued matrix Cauchy transform.
For invertible $G$ near $G_0$ set
\[
 \mathcal E(w,G):=G^{-1}-w1_{2m}+A+\Theta(G).
\]
Its derivative with respect to $G$ at $(0,G_0)$ is the invertible complex-linear map
\[
 L(X):=-G_0^{-1}(\id-\mathcal S)(X)G_0^{-1}.
\]
This produces a unique local analytic solution $g(w)$ with $g(0)=G_0$.
Indeed, on a sufficiently small closed ball about $G_0$ the map $G\mapsto G-L^{-1}\mathcal E(w,G)$ has derivative in $G$ of norm at most $1/2$.
For $|w|$ small it maps that ball into itself, because its value at the center tends to $G_0$.
Its iterates converge uniformly to the unique fixed point, and locally uniform convergence of the holomorphic iterates gives analyticity in $w$.

Adjoint symmetry and local uniqueness give $g(\overline w)=g(w)^*$.
Differentiating the equation at zero and using \eqref{eq:strict-full-stability} yields
\begin{equation}\label{eq:gap-derivative}
 g'(0)=-(\id-\mathcal S)^{-1}(G_0^2)=-D<0.
\end{equation}
Thus $\operatorname{Im}g(i\varepsilon)=-\varepsilon D+O(\varepsilon^2)<0$ for small $\varepsilon>0$.
Dyson uniqueness identifies $g(i\varepsilon)$ with the matrix Cauchy transform of $\mathscr H$.
The identity theorem extends this equality to the upper half of a small disk.
Since $g$ is selfadjoint on the real diameter, the spectral-gap criterion in Section~\ref{subsec:gap-tools} shows that $\mathscr H$ is boundedly invertible.
As $\mathscr H^2=\operatorname{diag}(SS^*,S^*S)$, this is equivalent to bounded invertibility of $S$.

\subsection{Attainment and the unique minimizer}
Conversely, suppose that $S$ is boundedly invertible.
The minimizing trial in \eqref{eq:scaled-minimum} is
\[
 b_\varepsilon=E_m((S^*S+\varepsilon^21_m)^{-1}).
\]
Norm continuity of inversion in a spectral gap gives
\[
 b_\varepsilon\longrightarrow b_*=E_m((S^*S)^{-1})>0,
 \qquad Y_\varepsilon\longrightarrow E_m((S^*)^{-1}).
\]
The matrices $b_\varepsilon$ stay bounded, so $\varepsilon^2\Tr b_\varepsilon\to0$.
Passing to the limit in \eqref{eq:scaled-minimum} and using continuity of $\mathcal V$ shows that $b_*$ attains its infimum.
Continuity of $\mathcal T_b^{-1}(a_0)$ gives $Z_{b_*}=E_m((S^*)^{-1})$.

Every attained minimum lies in the open positive cone and is a critical point of the smooth real function $\mathcal V$.
This also applies to the expression minimized in \eqref{eq:main-determinant}, which equals $\exp(\mathcal V(b)/(2m))$ for every $b>0$.
The preceding subsection has shown that every critical point forces invertibility.
It also determines that critical point uniquely: once $S$ is invertible, the continued matrix Cauchy transform satisfies
\[
 G_0=-E_{2m}(\mathscr H^{-1}),\qquad
 g'(0)=-E_{2m}(\mathscr H^{-2}).
\]
The upper right block of the first identity determines $Y=E_m((S^*)^{-1})$, and the lower diagonal block of \eqref{eq:gap-derivative} determines $b=E_m((S^*S)^{-1})$.
This proves all equivalences and the uniqueness assertion in Theorem~\ref{thm:main}.

\begin{remark}[Capacity and attainment]\label{rem:capacity-attainment}
Attainment of homogeneous capacity is related to exact operator scaling \cite[Theorem 4.7]{Gur04}, and is a different question from attainment in Theorem~\ref{thm:main}.
For $m=1$, $a_0=0$ and $\eta(b)=b$, the homogeneous quotient $\eta(b)/b$ is identically one, whereas $(1+\eta(b))/b=1+1/b$ reaches its infimum only as $b\to\infty$.
The corresponding standard semicircular element is not invertible.
\end{remark}
